\begin{explanationbox}
	\textbf{\textcolor{CExplanation}{一句话直觉}}

	椭圆上一点到左、右焦点的距离，等于半长轴 $a$ 加上或减去离心率 $e$ 与该点横坐标 $x_0$ 的乘积。

	\medskip
	\textbf{\textcolor{CExplanation}{核心拆解}}
	\begin{description}[style=nextline, leftmargin=2.8em, labelsep=0.5em, itemsep=0.45em]
		\item[\textbf{要点一（符号对应）：}] 左焦半径 $|PF_1|$ 取“$+$”号（$a+ex_0$），右焦半径 $|PF_2|$ 取“$-$”号（$a-ex_0$），可用“点向右、左距离变长”当场判定符号。
		\item[\textbf{要点二（范围约束）：}] 由于 $x_0\in[-a,a]$，可得焦半径的取值范围为 $[a-c, a+c]$，其中 $c=ae$。
	\end{description}

	\medskip
	\textbf{\textcolor{CExplanation}{几何本质（第二定义）}}

	由椭圆第二定义：点 $P$ 到右焦点 $F_2$ 的距离等于到右准线 $x=a^2/c$ 的距离乘以 $e$，即 $|PF_2| = e\cdot\left(\frac{a^2}{c} - x_0\right) = a - ex_0$；左焦点同理，$|PF_1| = e\cdot\left(x_0 + \frac{a^2}{c}\right) = a + ex_0$。

	\medskip
	\textbf{\textcolor{CExplanation}{代数意义（线性表达）}}

	将原本含平方的根式距离转化为横坐标的线性函数，大大简化计算，使距离与坐标直接对应。

	\medskip
	\textbf{\textcolor{CExplanation}{考点价值（常见题型）}}
	\begin{itemize}[leftmargin=*, itemsep=0.5em, topsep=0.4em]
		\item \textbf{考法一（直接代入）：} 已知 $P$ 的横坐标，直接代入公式求得 $|PF_1|,|PF_2|$。
		\item \textbf{考法二（最值问题）：} 利用 $x_0$ 的范围 $-a\le x_0\le a$ 确定焦半径的取值范围，或求相关量的最值。
	\end{itemize}

	\medskip
	\textbf{\textcolor{CExplanation}{顿悟点}}

	\textbf{焦半径公式是椭圆第一定义与第二定义的统一：由 $|PF_1|+|PF_2|=2a$ 结合第二定义可直接推出。}

	\medskip
	\textbf{\textcolor{CExplanation}{使用场景}}
	\begin{itemize}[leftmargin=*, itemsep=0.5em, topsep=0.4em]
		\item \textbf{场景一（已知坐标计算）：} 给出椭圆上点的坐标，求该点到焦点的距离。
		\item \textbf{场景二（综合最值问题）：} 在椭圆中涉及线段长度、和差最值时，用焦半径公式简化。
	\end{itemize}


\medskip
\textbf{\textcolor{CExplanation}{不靠死记的三点定位法}}
\begin{itemize}[leftmargin=*, itemsep=0.35em, topsep=0.35em]
 \item \textbf{在中间：} $x_0=0$，点 $P$ 在短轴端点处，两条焦半径同为 $a$。
 \item \textbf{向右走：} $x_0>0$，离左焦点相对更远、离右焦点相对更近，所以左边“加”，右边“减”。
 \item \textbf{到最右端：} $P(a,0)$，左、右焦半径分别是 $a+c$ 和 $a-c$，正好等于 $a+ea$ 和 $a-ea$。
\end{itemize}
\textbf{提醒：} 三个特殊位置帮助理解、验算符号；\emph{不能单凭三个位置就证明所有点都满足公式}。

\medskip
\textbf{\textcolor{CExplanation}{为什么想到“平方差”？}}
若题中同时出现“到左焦点的距离”和“到右焦点的距离”，不妨记为 $r_1,r_2$。
椭圆第一定义已直接给出 \emph{和}：$r_1+r_2=2a$。
因为两个焦点关于原点对称，写出两条距离的\emph{平方差}，共同的 $y_0^2$ 立即消失：
\[
 r_1^2-r_2^2=\big[(x_0+c)^2+y_0^2\big]
 -\big[(x_0-c)^2+y_0^2\big]=4cx_0.
\]
再除以已知的和 $2a$，就得到 $r_1-r_2=2ex_0$。
\textbf{这是从题目条件自然想到焦半径公式的路线：先看“和”，再求“差”，最后求出两条距离。}

\medskip
\textbf{\textcolor{CExplanation}{准线法的几何画面}}
椭圆右准线位于 $x=\frac{a^2}{c}$，在右顶点 $x=a$ 的更右侧；
因此从椭圆上任一点 $P$ 到右准线的水平距离为 $\frac{a^2}{c}-x_0$。
椭圆第二定义把它乘以 $e$，恰好得到右焦半径
$e(\frac{a^2}{c}-x_0)=a-ex_0$；左侧同理。
\textbf{距离对坐标的线性关系，来源于准线到点的水平距离。}

\medskip
\textbf{\textcolor{CExplanation}{顺手得到的最值}}
\[
 |PF_1|\,|PF_2|=(a+ex_0)(a-ex_0)=a^2-e^2x_0^2.
\]
因为 $-a\le x_0\le a$，所以
\[
 \boxed{b^2\le |PF_1|\,|PF_2|\le a^2}.
\]
乘积在 $x_0=0$ 时最大、在 $x_0=\pm a$ 时最小。
\end{explanationbox}
